\section{Revisión manuscrita y formas aceptadas}\label{sec:revision}

\subsection{Mis primeras observaciones}
Al probar los signos por separado, me resultaron fáciles de recordar y entender.
Sin embargo, Colección y Fuente me costaban más al escribirlos.

Después de probar las formas simplificadas B, las encontré más fáciles de
escribir y de distinguir del resto. Por eso las adopté como las formas habituales
de \textbf{Colección K06} y \textbf{Fuente K16}. Esta valoración recoge mi
experiencia al escribir; todavía no tengo mediciones de velocidad o recuerdo.

\subsection{Forma anterior y forma habitual}
Conservo las formas A para reconocer documentos anteriores. Los alias
V06/V16 siguen dibujando B; H06/H16 dibujan A. Estos códigos auxiliares no crean
significados nuevos. En la escritura corriente se usan K06 y K16.
\begin{center}
\begin{tabular}{@{}lcc@{}}
\toprule & A: anterior & B: habitual \\ \midrule
Colección & \KalaSimbolo[23mm]{H06} & \KalaSimbolo[23mm]{K06} \\[4mm]
Fuente & \KalaSimbolo[23mm]{H16} & \KalaSimbolo[23mm]{K16} \\ \bottomrule
\end{tabular}
\end{center}

\textbf{Colección.} Arco abierto y zigzag continuo; dos trayectos de dibujo.
Orden sugerido: arco de arriba abajo; después el zigzag desde su extremo superior
izquierdo. La forma anterior necesitaba cuatro trayectos.

\textbf{Fuente.} Silueta de documento con esquina inclinada y barra interior;
dos trayectos. Orden sugerido: contorno desde la esquina inferior izquierda,
subiendo y rodeando la figura; después la barra. La forma anterior tenía tres.
El contorno pertenece al glifo, no es un grupo sintáctico.

\subsection{Encuentros de líneas corregidos}
Al revisar las imágenes detecté dos encuentros de líneas que necesitaban ajuste:
\begin{itemize}
\item \textbf{Mientras K22:} un segmento descendente se superponía parcialmente al
arco. Eliminé ese segmento y uní arco y barra en un único trayecto continuo.
\item \textbf{Espacio K08:} la barra interior sobrepasaba el lado oblicuo. Ahora
termina sobre ese lado, formando una unión limpia.
\end{itemize}
\KalaMostrar{\KalaSimbolo[24mm]{K22}\qquad\KalaSimbolo[24mm]{K08}}
Las correcciones conservan los significados. Las superposiciones señaladas no
codificaban otra operación ni eran trazos que hubiera que repasar al escribir.
Otros signos sí tienen uniones o cruces definidos; estos deben verse con claridad.

\subsection{Práctica que sigue siendo útil}
Conservo la comparación A/B del cuaderno como parte de la evolución del diseño.
Ya elegí B. En el anexo de comprobaciones reúno la revisión técnica de T02/T03
y T04/T05. La hoja de lectura me permite anotar mi interpretación antes de
consultar las soluciones. Con T16 y T15 puedo practicar una composición y un programa.

También tengo preparado M01 para volver a leerlo otro día. Dejo las casillas
del cuaderno vacías hasta realizar cada ejercicio.
